\rhead{DS501: Predictive Modeling \& Analytics}
\phantomsection 
\section*{Predictive Modeling \& Analytics (DS501)}
\label{major:DS_5}

\noindent \small{\hyperref[main:scheme]{$\leftarrow$ Back to Scheme}} \vspace{0.5cm}

\noindent \textbf{Credits:} 3 (2-0-2)

\subsection*{Theory}
\noindent\textbf{Unit 1} \\
\textit{Predictive Modeling Foundations:} Supervised vs. unsupervised learning review, Regression vs. classification frameworks, Model evaluation metrics (MAE, RMSE, R² for regression; precision, recall, F1, AUC for classification), Cross-validation strategies (k-fold, stratified, time-series), Bias-variance tradeoff revisited, Overfitting detection (learning curves, validation monitoring), Model selection criteria (AIC, BIC, cross-validated scores).

\vspace{0.3cm}

\noindent\textbf{Unit 2} \\
\textit{Ensemble Methods and Bagging:} Bagging fundamentals and variance reduction, Random Forests (feature bagging, bootstrap aggregating), Extra Trees and feature importance interpretation, Out-of-bag error estimation, Gradient Boosting Machines (GBM) - forward stage-wise additive modeling, XGBoost architecture (regularized objectives, second-order gradients), LightGBM and CatBoost optimizations.

\vspace{0.3cm}

\noindent\textbf{Unit 3} \\
\textit{Time Series Forecasting:} Time series components (trend, seasonality, cyclicity, irregularity), Stationarity testing (ADF, KPSS tests), ARIMA/SARIMA models (ACF, PACF analysis, differencing), Exponential smoothing methods (Holt-Winters), Prophet forecasting framework, LSTM/GRU for sequential forecasting, Cross-validation for time series (purged/time-series splits), Anomaly detection in temporal data.

\vspace{0.3cm}

\noindent\textbf{Unit 4} \\
\textit{Survival Analysis and Imbalanced Learning:} Survival function, hazard rates, Kaplan-Meier estimator, Cox proportional hazards model, Time-to-event prediction metrics (C-index, Brier score), Imbalanced classification techniques (SMOTE, ADASYN, undersampling), Cost-sensitive learning, Anomaly detection algorithms (isolation forest, one-class SVM, autoencoders), Class imbalance evaluation (PR-AUC, Matthews correlation coefficient).

\vspace{0.3cm}

\noindent\textbf{Unit 5} \\
\textit{Model Deployment, Monitoring, and MLOps:} Model serialization (joblib, ONNX), Containerization (Docker for ML models), API serving (Flask, FastAPI, TensorFlow Serving), Model monitoring (drift detection, performance degradation), A/B testing and champion/challenger patterns, Automated retraining pipelines, Explainable AI techniques (SHAP, LIME, partial dependence), Production ML challenges (data versioning, lineage tracking).

\newpage

\subsection*{Practical}
\noindent\textbf{Lab 1: Model Evaluation and Validation} \\
Focuses on implementing comprehensive cross-validation strategies and model selection pipelines.

\vspace{0.2cm}

\noindent\textbf{Lab 2: Ensemble Learning Implementation} \\
Focuses on Random Forest and XGBoost training, hyperparameter tuning, and feature importance analysis.

\vspace{0.2cm}

\noindent\textbf{Lab 3: ARIMA and Exponential Smoothing} \\
Focuses on time series decomposition, stationarity testing, and classical forecasting models.

\vspace{0.2cm}

\noindent\textbf{Lab 4: Deep Learning for Time Series} \\
Focuses on LSTM/GRU forecasting with attention mechanisms and multivariate inputs.

\vspace{0.2cm}

\noindent\textbf{Lab 5: Survival Analysis Implementation} \\
Focuses on Kaplan-Meier curves, Cox models, and time-to-event prediction evaluation.

\vspace{0.2cm}

\noindent\textbf{Lab 6: Imbalanced Classification Techniques} \\
Focuses on SMOTE oversampling, cost-sensitive learning, and PR-AUC optimization.

\vspace{0.2cm}

\noindent\textbf{Lab 7: Anomaly Detection Systems} \\
Focuses on isolation forest, autoencoders, and multivariate anomaly detection pipelines.

\vspace{0.2cm}

\noindent\textbf{Lab 8: Model Deployment APIs} \\
Focuses on containerized model serving with FastAPI and performance benchmarking.

\vspace{0.2cm}

\noindent\textbf{Lab 9: MLOps Pipeline Development} \\
Focuses on automated retraining, drift detection, and A/B testing implementation.

\vspace{0.2cm}

\noindent\textbf{Lab 10: Production Analytics Project} \\
Focuses on complete predictive analytics solution from data ingestion to stakeholder dashboards.

\vspace{0.2cm}

\newpage
\rhead{Curriculum Schema}
